% Generated draft: requires mathematical and source review.
\documentclass[11pt]{article}
\usepackage[margin=25mm]{geometry}
\usepackage{fontspec}
\setmainfont{Latin Modern Roman}
\usepackage{amsmath,amssymb,mathtools,unicode-math}
\setmathfont{Latin Modern Math}
\usepackage{enumitem,needspace,xurl,hyperref}
\setlength{\parindent}{0pt}
\setlength{\parskip}{5pt}
\setlength{\emergencystretch}{3em}
\newcommand{\N}{\mathbb{N}}
\newcommand{\Z}{\mathbb{Z}}
\newcommand{\Q}{\mathbb{Q}}
\newcommand{\R}{\mathbb{R}}
\newcommand{\C}{\mathbb{C}}
\newcommand{\F}{\mathbb{F}}
\newcommand{\A}{\mathbb{A}}
\newcommand{\PP}{\mathbb{P}}
\newcommand{\E}{\mathbb{E}}
\newcommand{\eps}{\varepsilon}
\newcommand{\dd}{\,\mathrm{d}}
\begin{document}
\section*{UnsolvedMath ID: 30002530}
\textbf{Source problem code:} catalogue entry
\subsection{Definitions and mathematical statement}
A group \(G\) is torsion-free if, for every \(g\in G\) and every integer \(n\geq 1\), the equality \(g^n=1\) implies \(g=1\). A K\"ahler group is a group isomorphic to \(\pi_1(X,x)\) for some connected compact K\"ahler manifold \(X\) and base point \(x\in X\). Here a K\"ahler manifold is a complex manifold carrying a Hermitian metric whose associated real two-form is closed.

Using the standard complex-geometric interpretation of the source, the question is whether every torsion-free K\"ahler group \(G\) admits a realization
\[
G\cong\pi_1(Y(\C),y)
\]
for some smooth connected projective variety \(Y\) over \(\C\) and some \(y\in Y(\C)\), with the fundamental group taken in the analytic topology. Equivalently, is
\[
\forall G\quad
\bigl(G\text{ is a torsion-free K\"ahler group}\bigr)
\Longrightarrow
\bigl(\exists\,Y/\C\text{ smooth, connected and projective},\ \exists y\in Y(\C):\ G\cong\pi_1(Y(\C),y)\bigr)
\]
true? No equality of dimensions between a K\"ahler realization and a projective realization is required.

The supplied record labels the problem partially solved, while its accompanying assessment describes this full projective-realization question as remaining open and mentions a substantial positive result for linear representations. The precise scope and hypotheses of that representation-theoretic result are not supplied; it is not incorporated as a theorem or as a restriction on the question above.
\subsection{Short English statement}
Does every torsion-free group that occurs as the fundamental group of a connected compact K\"ahler manifold also occur as the fundamental group of a smooth connected complex projective variety? The supplied dated assessment describes the full question as open and reports related progress for linear representations.
\subsection{Sources}
\begin{itemize}
\item Supplied local registry record: \texttt{data/unsolvedmath/problems.json}, permanent UnsolvedMath ID \(30002530\), source problem code \texttt{catalogue entry}, title ``Projective Realization of Torsion-Free K\"ahler Groups.'' The record gives the proposed year as 2014 and identifies no individual proposer.
\item Source report: \emph{Mini-Workshop: K\"ahler Groups} (2014), Oberwolfach Reports, pp.\ 531--553, DOI \texttt{10.4171/owr/2014/09}. \url{https://doi.org/10.4171/owr/2014/09}. The supplied background describes this collection as drawn from workshop reports and their open-problem discussions.
\item Supplied literature reference: Campana, Claudon, and Eyssidieux, 2014. \url{https://www.numdam.org/item/10.5802/jep.12.pdf}. The supplied assessment attributes a substantial positive result for linear representations to these authors, without stating its precise hypotheses or conclusion.
\item Supplied literature reference: Oberwolfach Report 9, 2014. \url{https://doi.org/10.4171/owr/2014/09}.
\item Literature assessment date: \texttt{2026-08-21}. This is supplied dated curation derived from the supplied OWR curation, not an independently verified literature assessment. It labels the record partially solved while describing the full torsion-free K\"ahler group projectivity question as remaining open; time-sensitive status and cited claims require independent verification.
\item The supplied statement status is \texttt{corrected\_verified}; its accompanying explanation says that mathematical content was preserved while notation, typography, or omitted definitions were repaired for publication. This is a supplied source-verification qualification, not verification performed for this draft.
\end{itemize}
\end{document}
